\section{Forcing rank by split extensions}\label{sec:rank}

We begin with the algebraic part. It uses only a free subgroup with a
specified basis. The geometry will enter when we construct a group
containing many such subgroups and a conflicting rank inequality.

\subsection{A single attachment}

Let \(F\) be free on \(e_{uv}\), \(1\le u,v\le k\), and define
\[
 d:F\longrightarrow\Z,\qquad d(e_{uv})=\delta_{uv}.
\]
Let \(U=\langle s,\gamma\mid[s,\gamma]=1\rangle\cong\Z^2\).
The formulas
\[
 hsh^{-1}=s,\qquad h\gamma h^{-1}=s^{-d(h)}\gamma
 \quad(h\in F)
\]
define an action of \(F\) on \(U\): in additive coordinates the action
is the integral shear \((a,b)\mapsto(a-d(h)b,b)\).
Set
\begin{equation}\label{eq:attachment}
 A=U\rtimes F.
\end{equation}
Both factors embed, \(s\) is central of infinite order, and the projection
\(A\to F\) is a retraction with kernel \(U\).
Rearranging the action gives
\begin{equation}\label{eq:shear}
 \gamma h\gamma^{-1}=s^{d(h)}h.
\end{equation}
In particular a finite presentation of \(A\) is
\begin{equation}\label{eq:attachment-presentation}
 \left\langle e_{uv},s,\gamma\ \middle|\
 [s,\gamma]=1,\ [s,e_{uv}]=1,\
 \gamma e_{uv}\gamma^{-1}=s^{\delta_{uv}}e_{uv}
 \right\rangle.
\end{equation}

\begin{lemma}[Rank forcing]\label{lem:rank-forcing}
For every field \(K\) and every injective homomorphism \(K[A]\to D\)
to a division ring, the matrix
\[
 M=(e_{uv}-1)_{u,v=1}^{k}
\]
is invertible over \(D\).
\end{lemma}

\begin{proof}
Identify \(K[A]\) with its image and let \(E\subseteq D\) be its
division closure. Since \(s\) commutes with \(K[A]\), it commutes with
every element obtained by ring operations and inverses. Thus
\(s\in Z(E)\). Its powers are pairwise distinct by faithfulness and
the infinite order of \(s\) in \(A\).

Put
\[
 \Delta=\diag(e_{11},\ldots,e_{kk}),\qquad
 B=M-\Delta,\qquad C=\Delta^{-1}B.
\]
The diagonal entries of \(B\) are \(-1\). Off the diagonal,
\(d(e_{uv})=0\), so conjugation by \(\gamma\) fixes \(B\).
Equation~\eqref{eq:shear} therefore gives,
for every \(n\ge0\),
\begin{equation}\label{eq:conjugate-core}
 (\gamma^n I)M(\gamma^{-n}I)
 =s^n\Delta+B=\Delta(s^nI+C).
\end{equation}
Here \(\Delta\) is invertible, and all entries lie in \(E\).

Suppose \(M\) were singular over \(E\).
For \(n=0,\ldots,k\), equation~\eqref{eq:conjugate-core} gives a nonzero
vector \(v_n\in E^k\) with
\[
 Cv_n=v_n(-s^n).
\]
Eigenvectors for distinct central eigenvalues are right-linearly
independent. Indeed, take a nonzero dependence with the smallest
possible number of terms. Apply \(C\), then subtract the original
dependence multiplied on the right by one of the eigenvalues
\(\lambda_0\). Centrality gives a shorter dependence
\[
 \sum_{\lambda_j\ne\lambda_0}v_j a_j(\lambda_j-\lambda_0)=0
\]
in which all coefficients remain nonzero. This is impossible.
The \(k+1\) vectors \(v_0,\ldots,v_k\) cannot be independent in \(E^k\).
Thus \(M\) is invertible over \(E\), and its inverse also works over \(D\).
\end{proof}

Only the division closure of the attachment is used in this proof.
The element \(s\) need not be central in a larger division ring
generated by other parts of the construction.

\subsection{The retraction and its kernel}

Let \(G\) contain finitely many subgroups \(F_i\), each free on an
array of the form above. Construct \(A_i=U_i\rtimes F_i\) separately
for each \(i\). Let \(\Gamma\) be the fundamental group of the star
of groups with central vertex \(G\), leaf vertices \(A_i\), and edge
groups \(F_i\):
\begin{equation}\label{eq:star}
 \Gamma=G*_{F_i,\ i\in I}A_i.
\end{equation}
This notation means successive amalgamated free products along the
specified copies of \(F_i\), without imposing any further relations.
All edge maps are injective. The normal form theorem for amalgams
embeds \(G\) and every \(A_i\) in \(\Gamma\); see~\cite[Theorem~11 and Corollary~3 in I.5.2]{SerreTrees}.

The identity on \(G\) and the projections \(A_i\to F_i\le G\)
agree on the edge groups. They define a retraction
\(\rho:\Gamma\to G\).

\begin{lemma}\label{lem:kernel}
There is a free-product decomposition
\begin{equation}\label{eq:kernel}
 \ker\rho\cong
 *_{i\in I,\ gF_i\in G/F_i}gU_i g^{-1}.
\end{equation}
In particular, the kernel is a free product of copies of \(\Z^2\).
\end{lemma}

\begin{proof}
Write \(L=\ker\rho\) and restrict the Bass--Serre tree action to \(L\).
The stabilizer in \(L\) of an edge or a central-type vertex is trivial,
because \(\rho\) is injective on every conjugate of \(F_i\) and of \(G\).
A leaf stabilizer is a conjugate of
\(L\cap A_i=U_i\).

The quotient graph can be computed exactly. For a vertex or edge group
\(B\), the map induced by \(\rho\) identifies
\[
 L\backslash\Gamma/B\ \cong\ G/\rho(B).
\]
There is one central vertex. For every \(i\), the leaf vertices and
edges are both indexed by \(G/F_i\), and each edge joins the central
vertex to its corresponding leaf. The quotient graph is a star, hence
a tree. Its edge groups are trivial and its leaf groups are the
displayed conjugates of \(U_i\). The fundamental group of this tree of
groups is their free product~\cite[Theorem~13 in I.5.4]{SerreTrees}.
\end{proof}

\subsection{Unique product and diffuseness}\label{sec:diffuse}

An element \(a\) of a finite subset \(S\) of a group is \emph{extremal}
if, for each \(h\ne1\), at least one of \(ha,h^{-1}a\) lies outside \(S\).
A group is \emph{diffuse} if every finite subset with at least two
elements has at least two extremal elements. Diffuse groups have unique
product, and free products of diffuse groups are diffuse
\cite[Proposition~1.1 and Theorem~1.2(2)]{Bowditch}.
These statements apply to arbitrary free products, since a finite
subset uses only finitely many factors. In particular, the kernel
in~\eqref{eq:kernel} is diffuse: each factor \(\Z^2\) is ordered
lexicographically, and the endpoints of a finite subset are extremal.

We record the extension arguments explicitly; the diffuse case is
\cite[Theorem~1.2(1)]{Bowditch}.

\begin{lemma}\label{lem:extension}
If \(1\to L\to H\overset{\rho}{\to}Q\to1\) is exact and \(L,Q\)
have unique product, then \(H\) has unique product. If \(L,Q\) are
diffuse, then \(H\) is diffuse.
\end{lemma}

\begin{proof}
For unique product, take nonempty finite \(P,R\subset H\). Choose
a uniquely represented product \(ab\) in \(\rho(P)\rho(R)\).
All representations of an element over \(ab\) must use
\[
 P_a=P\cap\rho^{-1}(a),\qquad R_b=R\cap\rho^{-1}(b).
\]
Choose \(p_0\in P_a\), \(r_0\in R_b\), and write
\(P_a=p_0P'\), \(R_b=R'r_0\), with \(P',R'\subset L\).
A uniquely represented element of \(P'R'\) gives one of \(PR\).

For diffuseness, take a finite \(S\subset H\) with \(|S|\ge2\).
If \(\rho(S)\) has at least two elements, choose two extremal elements
of \(\rho(S)\). In each of their fibers choose an element extremal for
the action of \(L\), using diffuseness of \(L\), or the only element
if the fiber is a singleton. Each selected element is extremal in
\(S\): an element outside \(L\) is controlled by the quotient, and an
element in \(L\) is controlled by the fiber.
If \(\rho(S)\) is a singleton, identify that fiber with an \(L\)-coset
and take two extremal elements there. They are again extremal in \(S\).
\end{proof}

For completeness, unique product gives domainhood directly. If
\(\alpha=\sum_{a\in P}\alpha_a a\) and
\(\beta=\sum_{b\in R}\beta_b b\) have nonzero coefficients on their
finite supports, choose a uniquely represented \(ab\in PR\).
Its coefficient in \(\alpha\beta\) is
\(\alpha_a\beta_b\ne0\).

\begin{proof}[Proof of Theorem~\ref{thm:principle}]
For each \(i\), attach~\eqref{eq:attachment} to the specified free
basis of \(F_i\), with \(k=r_i-1\), and form~\eqref{eq:star}.
Lemmas~\ref{lem:kernel} and~\ref{lem:extension} give unique product,
and give diffuseness when \(G\) is diffuse.

Suppose \(K[\Gamma]\) embeds in a division ring \(D\). Its restriction
to each \(K[A_i]\) is injective.
Lemma~\ref{lem:rank-forcing} makes every matrix
\((e_{i,uv}-1)\) invertible. Subtract the first row of the normalized
\(Y_i\) from each of its other rows. The resulting upper block
triangular matrix has diagonal blocks \(1\) and
\((e_{i,uv}-1)\). Thus \(\rk_D Y_i=r_i\) for every \(i\), contrary to
the assumed rank inequality.

If \(G\) is finitely presented, add two generators \(s_i,\gamma_i\)
for each \(i\), together with the finitely many relations
in~\eqref{eq:attachment-presentation}, using the specified words for
the \(e_{i,uv}\) in \(G\). This is a finite presentation of \(\Gamma\).
\end{proof}
